\scanpage{105}
% The heading is enclosed in a hand-drawn frame.
\uline{Poids de la cohomologie \inserted{$l$-adique} d'une variété singulière ou/et non propre \inserted{(sur corps $k$ parfait)}.}

\textcircled{A} \uline{Coefficients constants}

\textcircled{1} Soit d'abord $X$ lisse et quasi-projective, donc par résolution forte (que nous admettons) on aura
\[
 X=X'-Y \qquad Y=\bigcup_i Y_i
\]
$X'$ lisse et projective, $Y_i$ diviseurs lisses \inserted{irréd.} à croisements normaux mutuels. \deleted{\unsure{Soit}} On a donc
\[
\begin{aligned}
\inserted{H^{i-1}(X')\to H^{i-1}_{\unsure{!|{}}}(Y)}&\\[-3pt]
\deleted{\unsure{R^if_!}}&\to H^i_!(X)\to H^i(X')\to H^i(Y)\to H^{i+1}_!(X)\to\ldots
\end{aligned}
\]

% Small diagram in the left margin.
\diagram{\begin{tikzcd}[ampersand replacement=\&,row sep=small,column sep=small]
X\arrow[r,phantom,"\subset"]\arrow[d,"f"'] \& X'\arrow[dl,"f'"]\\
\operatorname{Sp}\unsure{_{\mathrm e}|{}|\mathrm{ec}}\,k \&
\end{tikzcd}}

\deleted{\unsure{donc}} donc les poids dans $H^i_!(X)$ (qui se déterminent via val. absolues de val. propres de Frob, par réduction au cas d'un corps de type fini \ldots) est contenu dans ceux de $H^i(X')$ [réduits à $i$ par Weil] et ceux de $H^{i-1}(Y)$.

Or on
\[
 H^*(Y)\Longleftarrow E_2^{pq}
 =H^p\bigl(\unsure{i_0\ldots i_r|i_0\ldots i_p}\longmapsto H^q(Y_{i_0}\cap\ldots\cap Y_{i_{\unsure{r|p}}})\bigr)
 \left\{\begin{aligned}
 &\text{de poids }q\\
 &\text{de niveau }\le q\\
 &[\text{intervient si }\\[-2pt]
 &i=p+q)\le n-1]\\
 &\text{de niveau }\le\unsure{2n-i-q}\\
 &\inserted{\unsure{-q|{-}p}}\\
 &[\unsure{\text{et}|\text{si}}\, i=p+q\ge n]
 \end{aligned}\right.
\]
\margin{$(2n-i-1)-q$}
\margin{(on peut se borner ici aux cochaînes \uline{alternées}}

\deleted{Donc} Or ici $Y_{i_0}\cap\ldots\cap Y_{i_p}$ est projectif et lisse (de dimension $n-p-1$ si $X$ \deleted{\unsure{connexe}} purement de dim $n$).

Donc le poids du terme $E^{pq}_2$ est $q$, et on trouve que le poids de $H^j(Y)$ est $\le j$. Donc :

\deleted{\uline{\unsure{Lemme 1|Lemme}}} \inserted{\uline{Proposition 1}} $H^i_!(X)$ est de poids \inserted{entre $0$ et} $\le i$. La partie de poids $i$ (resp. de poids $\le i-1$) est le quotient image de $H^i_!(X)\to H^i(X')$ [resp l'image de $H^{i-1}(Y)\to H^i_!(X))$. De plus, si $i\ge n=\dim X$, alors le \uline{niveau} de $H^i_!(X)$ est $\le 2n-i$, et les \unsure{deuxièmes|derniers|deux} morceaux précédents sont de niveau $\le 2n-i-1$ ; de plus, il n'y a de poids qu'entre $2(i-n)$ [car il faut écrire $q\le 2(n-p-1)$ !] et $i$, i-e sur un intervalle d'amplitude $2n-i$ !\ldots

\margin{\uline{NB} Tous les morceaux de la filtration sont \textquotedblleft effectifs\textquotedblright}
\margin{En tous cas, niveau $\le\operatorname{Inf}(i,n)$}

Pour ce complément, notons que le niveau de $E^{pq}_2$ est $\le$ celui de $H^q(\underbrace{Y_{i_0}\cap\ldots\cap Y_{i_p}}_{\dim n-p-1})$, donc est
\[
\begin{aligned}
 {}\le{}&\deleted{\unsure{\operatorname{Inf}(q,}}\,2(\deleted{\unsure{\dim}}n-p-1)-q\,\unsure{\bigr)|\}|{}}\\
 ={}&\deleted{\unsure{\operatorname{Inf}(q,}}(2n-i-2)-p
 \inserted{\deleted{\unsure{2q}}=2(n-1)-i-p}
 \quad(\inserted{\text{si }}i=p+q)\ \text{donc }\le 2(n-1)-i
\end{aligned}
\]
% archive sheet number: 104
